

\chapter{BACKGROUND STUDY AND LITERATURE REVIEW}
\ChapterSection{Background Study}

\ChapterSubsection{Crop Disease Diagnosis as a Decision-Support Problem}

The Crop disease diagnosis in the field includes more than just the image classification. Farmers need immediate intervention strategies like specific chemicals that can be sprayed, dosage and how fast it should be sprayed, which the computer vision model cannot provide at their current state. Despite a lot of higher accuracy in the identification of plant pathogens in computer vision models, researchers rarely try to map these biological labels to practical agriculture practices. This gives rise to an application gap between the prediction done by algorithms and the decision support required from them. Nepal's agriculture accounts for a quarter of the national GDP, and most rural households depend on it\parencite{moaldAgricultureStats2080,nsoAgricultureCensus2021}.
Smallholder farmers of Nepal operate under strict budget constraints. And if the farmer selects the wrong pesticide or delays a response to a fungal disease can lead to loss of income \parencite{pandeyCropLossNepal2022}. 

Since the farmer's losses can be minimised using digital software, the software would be deployed depending on local infrastructure. From ICT analysis and studies by FAO, even as the number of farmers with smartphones grows, network connectivity is segmented in rural areas. Another issue would be a complicated user interface and software that only consists of English text, which make the farmers give up using agricultural applications eventually (Sharma, Singh, Rana and Kumar, 2017). Though having higher model accuracy in the test environment does not solve anything if deployment is limited due to local infrastructure. It would have practical implications if the diagnosis could be performed offline and local instructions written in Nepali, as well treatments, could be provided directly as steps instead of only the name of the disease \parencite{silvestriSmallholderTechAdoption2021}.

\ChapterSection{Literature Review}

\ChapterSubsection{Deep Learning for Plant Disease Detection}

\subsubsection{Dataset evolution and the lab-to-field gap}

Deep learning for plant pathology has increased with the PlantVillage study. Convolutional neural networks achieved over 99 percent accuracy in classifying 26 diseases with across 14 crop species \parencite{mohanty2016plantvillage}. Even though these results looked like a breakthrough. The images taken by the user can be varied based on the user device, environment, lighting and other factors, which would be different from the dataset that was used for the training; therefore, the model would likly the collapsed on these outlier data \parencite{singhPlantDoc2020,thakurPlantDiseaseFieldGap2022}.

The researchers have collected ``in-the-wild'' disease images from different sources online, and are trying to prove that standard transfer-learning models struggle outside the experimental environment \parencite{singhPlantDoc2020}.  Best-in-class machine learning models reach only 60 to 75 percent. This huge performance drop, often exceeding 25 percentage points, now opens up agricultural machine learning research \parencite{thakurPlantDiseaseFieldGap2022,li2021plantDiseaseReviewCNN}.

The machine learning models evolved from AlexNet and VGGNet through ResNet \parencite{heResNet2016}, to current models such as EfficientNet architectures that rely heavily on ImageNet pre-training \parencite{saleem2019plantDiseaseML,liu2021plantDiseaseReview}. Despite this, choosing the best-performing model means very little, and the real performance depends on the quality of the dataset and how the classes are balanced, including the specific augmentation tricks used during training. And testing these models on the original PlantVillage dataset inflates accuracy therefore, the true evaluations require messy, field-collected data \parencite{thakurPlantDiseaseFieldGap2022}.


\subsubsection{Existing mobile disease-detection systems}

Different mobile applications have tried to deliver deep-learning-based plant disease diagnosis
to farmers. Plantix, developed by PEAT, uses a convolutional neural network to classify disease,
nutrient deficiency, and pest damage from photographs, with treatment advice delivered in
multiple languages \parencite{petrellaPlantix2023}. PlantVillage Nuru, developed by Penn State
University in partnership with IITA and the FAO, provides on-device inference using TensorFlow
Lite and was specifically designed for offline use a critical feature for farmers in remote
areas with unreliable connectivity \parencite{mrishoPlantVillageNuru2020}.

Table~\ref{tab:existing-mobile-apps} compares published mobile plant-health applications along
dimensions relevant to Krishi Vaidya's design goals.

\begin{table}[!htbp]
\centering
\caption{Comparison of mobile plant-disease advisory systems}
\label{tab:existing-mobile-apps}
\small
\renewcommand{\arraystretch}{1.3}
\begin{tabularx}{\textwidth}{@{} 
  >{\raggedright\arraybackslash}p{2.8cm} 
  >{\raggedright\arraybackslash}X 
  >{\raggedright\arraybackslash}X 
  >{\raggedright\arraybackslash}X 
  >{\raggedright\arraybackslash}p{2cm} 
  >{\raggedright\arraybackslash}p{2cm} @{}}
\toprule
\textbf{System} & \textbf{Detection method} & \textbf{Offline support} & \textbf{Advisory output} & \textbf{Language} & \textbf{Crop scope} \\
\midrule
Plantix & CNN classification & Partial (capture offline, diagnose online) & Treatment plan, fertilizer info & 18+ languages & 30+ crops \\
PlantVillage Nuru & TFLite on-device & Full offline inference & Basic management advice & Swahili, English & Cassava, potato, others \\
eRice (IRRI) & Image upload + expert & No (server required) & Expert-validated advisory & English, local & Rice \\
Krishi Vaidya & YOLO + VLM + backend & Offline queue with cached advisory & Structured bilingual advisory & Nepali, English & 6 crops \\
\bottomrule
\end{tabularx}
\end{table}

Existing systems generally address one or two of these capabilities: on-device detection,
structured advisory generation, offline aware, and crop-management record keeping. None
of the reviewed systems integrates all four within a single workflow that also includes
VLM-based reasoning. Krishi Vaidya's contribution lies in this integration rather than in
surpassing any individual component's accuracy---the research gap is architectural and
workflow-oriented, not primarily accuracy-oriented.

\ChapterSubsection{Object Detection in Agricultural Vision}

\subsubsection{The YOLO detector family}

Object detection can provide the results of \emph{where} a plant structure exists, not only \emph{what} it is. The images taken by the farmers often contain poor lighting, low quality, and many degradation factors. The disease symptoms usually hide in small or overlapping fragments of a larger image of the plant. A good model would look for the structural elements like panicles and fruits, including torn or partially obscured leaves. And isolating these regions prevents the downstream diagnostic models from getting lost in the background noise of the normal or healthy portion.

The YOLO architecture includes the real-time agricultural space \parencite{jiangYOLOSurvey2022}. It can run faster, and it stays highly accurate. Ultralytics improved the performance further with YOLOv8, packing segmentation capabilities and pose estimation alongside standard detection, all wrapped in an extremely efficient training pipeline \parencite{ultralyticsYolov8}. Farm-based deployments of YOLO variants show mAP@50 scores in the range between 0.70 and 0.92 \parencite{tian2020AppleFruitYOLO,zhengSmallObjDetAgri2023}. This increased variance depends more on the specific crop profile and the volume of training data, alongside whether researchers relied on clean lab photos instead of messy field-plan disease pictures.

For Krishi Vaidya, we have choosed the version YOLOv8 nano (\texttt{yolov8n}). We chose this specific version of yolo because it can run on low-end devices. This version has only 3.2 million parameters, and it can run comfortably inside constrained mobile memory without relying on the cloud. The model finds the crop parts. It does not diagnose diseases. 

\subsubsection{Detection evaluation methodology}

The PASCAL VOC challenge set the original baseline. It evaluates object detection at a flat 0.50 Intersection-over-Union (IoU) threshold \parencite{everinghamPASCALVOC2010}. Then COCO updated the rules with benchmark demands absolute pixel precision. It averages scores across a range of thresholds from 0.50 all the way to 0.95, penalising the models that successfully identify objects but draw sloppy bounding boxes around them \parencite{linCOCO2014}.

It is important to understand that relying on a single metric hides fatal flaws, and also isolated numbers hide serious evaluation biases \parencite{padillaMetricsSurvey2021}. Krishi Vaidya tracks mAP alongside precision and recall, including the harmonic F1-score to uncover edge-case failures. 

Our architecture requires high recall. The YOLO model works as a first-stage scout for the Vision-Language Model (VLM) pipeline. Missed leaves forever lost. And the downstream diagnostic steps cannot analyse a crop region the model has ignored. Conversely, false positives rarely matter. If YOLO accidentally flags a patch of dirt as a leaf, the VLM will reject it later. Chapter 5 breaks down this exact threshold tradeoff.

\subsubsection{Small-object and dense-detection challenges}

Small crop fragments often break object detectors, and models struggle severely when targets clump together \parencite{zhengSmallObjDetAgri2023}. Krishi Vaidya classes include rice panicles alongside tightly packed grain clusters. These structures disappear visually. They could hide perfectly into the surrounding green background.

The low spatial resolution could omit the vital structural details. Human annotators for the dataset often struggle to draw precise bounding boxes around messy, overlapping stems, which introduces noisy ground-truth labels directly into the training dataset.

The YOLO results from our experiment show a similar conclusion. The model often misses rice panicles and brassica leaves. And this significant drop proves that the extreme case of small-object geometry rather than an inherent flaw in the training process.

\subsubsection{Quantization and deployment}


Trying to run the neural network on a low-end smartphone comes up the comprmised in compromised performance. The accuracy should be sacrificed against inference speed. Post-training quantisation reduces the model parameters. It squeezes full-precision floating-point (FP32) weights down into smaller formats like FP16 or INT8, reducing the file size, exchanging it with drops in performance \parencite{garg2022quantizationSurvey}. Krishi Vaidya relies on native Ultralytics export pipelines to build deployment variants \parencite{ultralyticsExport}. 

 Reducing the model down to INT8 improves detection for a few specific diseases while obliterating accuracy for other classes of diseases. Aggregate baseline scores mask these isolated collapses. Safer field deployment often requires tracking exactly how each category survives quantisation \parencite{paleyes2022MLSystemsChallenges}.
 
\ChapterSubsection{Vision-Language Models for Structured Diagnosis}

\subsubsection{From classification to structured reasoning}

A modern VLM takes a leaf image and returns a fully analyzed report. Not only is the exact disease printed and a confidence percentage predicted, a full list of associated symptom descriptions and the immediate next steps for the treatment are too. LLaVA set out the precedent and instructions that form the training backbone of all VLMs \parencite{liuLLaVA2023}. This has been redesigned past simple image captioning to understanding VQA and intricate information extraction \parencite{zhangVLMSurvey2024}. Agri-LLaVA was trained on fixed pests and diseases pair between over 200, across a 2-phase training protocol \parencite{wangAgriLLaVA2024}. 

The LLMI-CDP architecture has achieved similar results. Developers of LLMI-CDP architecture followed existing VLMs using Low-Rank Adaptation (LoRA) and domain-specific output formats to improve diagnostic accuracy \parencite{yangLLMCropDisease2024}. Krishi Vaidya ran LoRA adaptation over a pre-trained Qwen2.5-VL-3B-Instruct backbone. The training teaches the raw model to generate predictable, structured crop assessments.

\subsubsection{Parameter-efficient fine-tuning}

Updating a billion parameters of VLM requires compute that most student labs cannot afford. LoRA solves this hardware problem. It freezes the base weights and trains low-rank adapter matrices, injecting thes adapter matrices  straight into selected attention and feed-forward layers \parencite{hu2021lora}. The base model retains its complex learned representations without requiring a full parameter update. 

QLoRA pushes memory efficiency even further as it loads the frozen weights in 4-bit quantised form to fit massive architectures into standard GPU limits \parencite{dettmers2023qlora}. Krishi Vaidya uses both techniques. The Qwen2.5-VL base model runs under 4-bit NF4 quantisation while LoRA adapters train on the target modules. We completed all fine-tuning on a single 22 GB NVIDIA L4 GPU. 

With using these quantisation and adapter techniques made single-GPU training an actual possibility. Past transfer-learning research shows similar outcomes. Distribution shifts between strict training supervision and open inference environments remain a persistent hurdle \parencite{zhuangTransferLearningSurvey2021}.

\subsubsection{Structured output and parseability}

Natural-language text fails if the backend cannot extract the correct disease
name, confidence numbers, and specific instructions detailing
treatment steps. Unstructured responses break downstream storage.
\ChapterSubsection{Offline-First Mobile Architecture}

Agricultural mobile apps in developing regions face constant infrastructure failures. Field tools must function offline. \textcite{kleppmannLocalFirst2019} outlines the ``local-first'' software paradigm to solve this dependency. 

Engineers build such a system with using embedded databases. SQLite is used for the primary data storage.  Background processes handle synchronisation with the server if the mobile is connected to the internet \parencite{mehdiOfflineFirstMobile2022}. This setup reduces reliance on the cloud. 

Krishi Vaidya runs entirely on this approach. SQLite database on the mobile app holds the raw scan records alongside user settings, pending synchronisation queues with the server, and daily crop plots the farmer logged. Internet connections trigger a backend VLM diagnosis. If the user tries to scan the crop and there is no internet, then it will be saved into the offline queue and the request will be sent automatically upon internet connection. 
\ChapterSubsection{Backend Architecture and ML System Design}

 The real-world applications often require heavy transformation to turn raw model outputs into structured data. A server must authenticate incoming requests before directing the model service, parsing outputs, and updating the database. 

\textcite{sculleyMLTechDebt2015} presented the ``hidden technical debt'' of ML systems.  The major engineering effort goes into the surrounding infrastructure to manage data collection, build active serving setups, and handle complex configurations. The deployment of the end-to-end system complexity could extend deep into data pipelines and organisational limits \parencite{paleyes2022MLSystemsChallenges}. 

Krishi Vaidya runs on an Express and Mongoose API to manage this complexity. The implemented backend codebase isolates middleware functions from request controllers and raw database repositories. The uploaded images are validated thoroughly before it is send to controller layers. The backend then sends the request to the VLM service. Then it parses the raw text output before a secondary model call structures the final advisory against a standard schema. 

\ChapterSubsection{Research Gap and Positioning}

Current literature shows three clear and major gaps. Krishi Vaidya tackles them within the scope of a final-year project.

\textbf{Gap 1: Integration gap.} Most plant-disease systems can only do the disease classification, as they ignore the farmers' need to capture images and read structured advice. Basic tools like plot management or saving past scans are missing. There are some mobile apps such as Plantix and PlantVillage Nuru that provide some features that mentioned before \parencite{petrellaPlantix2023,mrishoPlantVillageNuru2020,mohanty2016plantvillage}. Despite they remain fragmented. A working diagnostic system must combine raw object detection with VLM interpretation and a backend built for offline delivery.

\textbf{Gap 2: Structured-output evaluation gap.} Agricultural VLM and LLM benchmarks report classification accuracy \parencite{wangAgriLLaVA2024,yangLLMCropDisease2024}. They often ignore whether a system can parse those outputs into usable data structures like JSON and others. Parseability is essential just as much as semantic correctness for a mobile application. Krishi Vaidya includes it as a core metric alongside standard accuracy.

\textbf{Gap 3: Contextual deployment gap.} Published benchmarks skip real-world deployment constraints \parencite{thakurPlantDiseaseFieldGap2022}. Most of them ignore offline operation and bilingual or multi-language support. Academic evaluations and tests miss how quantisation shifts model behaviour on a mobile phone linked to actual farm records. Krishi Vaidya evaluates these constraints upfront and end to end. 

Krishi Vaidya does not close these gaps. Rather, it is a student prototype. The major contribution of this project lies in the integration. The trained models, a validated backend, and an offline-aware app can form a cohesive diagnostic pipeline for the farmer with ease of access and usability.

